\section{Conclusion}
We presented a supervised method for adapting signal control across road networks by learning differences between feasible phase choices. A restricted reference predictor and an antisymmetric pairwise predictor are combined using target simulation data, then deployed without online refitting. The final controller selects source components or a target anchor using eight target development seeds before freezing its predictions and execution guard. On 56 paired seeds across two target networks, it reduces mean city-wise relative waiting by 2.42\% against MaxPressure, with a 95\% reduction interval of 0.66--4.58\% and 40 favorable seed means. It also lowers waiting relative to the tested target anchor and dense linear residual baseline. The PhasePressure point estimate is favorable, but its interval crosses zero. The evidence supports the complete source-selection procedure on these fixed networks and motivates capacity-matched tests of the source data's contribution on additional independent networks.
